\uuid{tgoL}
\titre{Racine}
\auteur{bodin ; hanani ; mzari}
\organisation{exo7}
\licence{CC BY-NC-SA 4.0}
\datecreate{2020-01-20}
\niveau{L1}
\module{Algèbre}
\chapitre{Polynôme, fraction rationnelle}
\sousChapitre{Racine, décomposition en facteurs irréductibles}
\difficulte{2}
\theme{Racines, factorisation, racine}
\exo7qcmid{lille-1.polynomes.021}

\contenu{
\begin{question}
\qclassification{105.03}
\qtags{motcle=racine}
Soit $P \in \Rr[X]$ un polynôme de degré $2n+1$ ($n\in\Nn^*$).
Quelles sont les affirmations vraies ?
\begin{answers}
    \good{$P$ peut admettre une racine complexe, qui ne soit pas réelle.}

    \good{$P$ admet au moins une racines réelle.}

    \bad{$P$ admet au moins deux racines réelles (comptées avec multiplicités).}

    \good{$P$ peut avoir $2n+1$ racines réelles distinctes.}
\end{answers}
\begin{explanations}
Il y a au plus $\deg P$ racines réelles (comptées avec multiplicité). Mais ici le degré est impair, donc $P$ admet au moins une racine réelle.
\end{explanations}
\end{question}
}
